% Exemplo: programação com um fluxograma em TikZ e um algoritmo com algorithm2e (palavras-chave
% em português pela opção portuguese).
% Compilar: latexmk -pdf programacao.tex          Testado com: algorithm2e 5.2, TeX Live 2024 ou mais recente
%%BEGIN:progpreambulo
\documentclass[12pt, a4paper]{article}
\usepackage[T1]{fontenc}
\usepackage[brazilian]{babel}
\usepackage{tikz}
\usetikzlibrary{shapes.geometric, shapes.misc, positioning}
\usepackage[portuguese, ruled, linesnumbered]{algorithm2e}
%%END:progpreambulo

\begin{document}
\pagestyle{empty}

%%BEGIN:progfluxograma
\begin{center}
\begin{tikzpicture}[node distance=7mm, font=\small,
    terminal/.style={draw, rounded rectangle, minimum width=2cm},
    processo/.style={draw, rectangle, minimum width=2.6cm, minimum height=7mm},
    decisao/.style={draw, diamond, aspect=2, inner sep=1pt},
    es/.style={draw, trapezium, trapezium left angle=70, trapezium right angle=110},
    seta/.style={-stealth}]
  \node[terminal]                     (ini)  {início};
  \node[es, below=of ini]             (le)   {ler $n_1$, $n_2$};
  \node[processo, below=of le]        (med)  {$m \gets (n_1 + n_2)/2$};
  \node[decisao, below=of med]        (dec)  {$m \ge 6$?};
  \node[es, below left=7mm and 10mm of dec]  (apr) {``aprovado''};
  \node[es, below right=7mm and 10mm of dec] (rep) {``reprovado''};
  \node[terminal, below=18mm of dec]  (fim)  {fim};
  \draw[seta] (ini) -- (le);
  \draw[seta] (le) -- (med);
  \draw[seta] (med) -- (dec);
  \draw[seta] (dec) -| node[pos=0.3, above] {sim} (apr);
  \draw[seta] (dec) -| node[pos=0.3, above] {não} (rep);
  \draw[seta] (apr) |- (fim);
  \draw[seta] (rep) |- (fim);
\end{tikzpicture}
\end{center}
%%END:progfluxograma

%%BEGIN:progalgoritmo
\begin{algorithm}[htbp]
  \caption{Fatorial de um número natural}
  \Entrada{um número natural $n$}
  \Saida{$n!$}
  $f \gets 1$\;
  \Para{$i \gets 2$ \Ate $n$}{
    $f \gets f \times i$\;
  }
  \Retorna{$f$}\;
\end{algorithm}
%%END:progalgoritmo

\end{document}
